\honest{We Change Our Minds Less Often Than We Think}{Eliezer Yudkowsky}{2007}

I open with Dale Griffin and Amos Tversky. They asked colleagues choosing between job offers how likely they were to pick each job. The average confidence in the predicted choice was ``a modest 66\%,'' but only 1 of the 24 chose the other job. \nb{In the paper this is an example of underconfidence. The authors explain it by confidence that follows the balance of arguments: a 2-to-1 balance of reasons feels like two-thirds certainty, although ``even a small advantage for job A over B would generally lead to the choice of A.''}

Reading this, on August 1st, 2003, at around 3 o'clock in the afternoon, changed the way I thought. Once I could guess my answer, I realized, I had ``in all probability, already decided.'' We change our minds less often than we think. \nb{The survey measured how well people predicted their choice, not when they decided. On the same page Griffin and Tversky note that people are sometimes overconfident in predicting their own behavior, and cite a study of students whose predictions of their own actions ``proved to be consistently overconfident.''} Most of the time we can guess our answer within half a second of hearing the question. I give no source for this.

That moment before we can guess is the only window for intelligence to act, ``In questions of choice, as in questions of fact.'' \nb{The survey concerned choices only.}

The principle of ``The Bottom Line'' is that only the real causes of your beliefs determine how well you do. You might think you could reach a belief by non-rational means, check its justification, and drop it if the check fails. But we change our minds ``much less often'' than we think. \nb{The survey concerned people weighing two job offers. It did not observe anyone checking a conclusion and failing to drop it.}

You can remember a time you changed your mind. Can you remember as easily all the times you did not? Between hindsight bias, fake causality, positive bias, anchoring and ``above all the dreaded confirmation bias,'' an idea that gets into your head is ``probably going to stay there.'' \nb{The finding closest to this, which the post does not mention, is belief perseverance. Ross, Lepper and Hubbard (1975) found that impressions persisted after the feedback behind them had been ``completely discredited.''}
